%% kapa: exam=1 round=33 year=2022 period=1 form=A subject=civil_law subject_ko=민법 seq=27 number=27
%% source: Q-Net 감정평가사 기출문제 (cst00309) artlSeq=5209997 file=raw/1008/5209997/2022년도 제33회 감정평가사 1차시험 1교시(민법,경제학,부동산학).pdf
%% conversion: kordoc 4.24.51 -> tools/gen1.py (auto); review=proofread-sonnet; flags=none
\begin{question}{27}
점유에 관한 설명으로 옳은 것은? (다툼이 있으면 판례에 따름)
\begin{choices}
  \item 미등기건물의 양수인은 그 건물에 관한 사실상의 처분권을 보유하더라도 건물부지를 점유하고 있다고 볼 수 없다.
  \item 건물 공유자 중 일부만이 당해 건물을 점유하고 있는 경우, 그 건물의 부지는 건물 공유자 전원이 공동으로 점유하는 것으로 볼 수 있다.
  \item 점유자의 권리적법추정 규정(민법 제200조)은 특별한 사정이 없는 한 등기된 부동산에도 적용된다.
  \item 선의의 점유자라도 본권에 관한 소에 패소한 때에는 그 패소판결이 확정된 때로부터 악의의 점유자로 본다.
  \item 진정한 소유자가 점유자를 상대로 소유권이전등기의 말소청구소송을 제기하여 점유자의 패소로 확정된 경우, 그 소가 제기된 때부터 점유자의 점유는 타주점유로 전환된다.
\end{choices}
\end{question}
\answer{2}
